% Part: sets-functions-relations
% Chapter: relations
% Section: orders

\documentclass[../../../include/open-logic-section]{subfiles}

\begin{document}

\olfileid{sfr}{rel}{ord}
\olsection{ક્રમો}

\begin{explain}
આપણી ઘણી સરખામણીઓમાં કેટલીક વસ્તુઓને, કોઈ ચોક્કસ બાબતમાં,
બીજી વસ્તુઓ “કરતાં નાની”, “સમાન” અથવા “કરતાં મોટી”
તરીકે વર્ણવવામાં આવે છે. તેમાં \emph{ક્રમ} સંબંધો આવે છે.
પણ ક્રમ સંબંધો જુદા જુદા પ્રકારના હોય છે. દાખલા તરીકે,
કેટલાકમાં કોઈ પણ બે વસ્તુઓ તુલનીય હોવી જરૂરી છે,
બીજામાં એવું નથી. કેટલાકમાં તાદાત્મ્યનો સમાવેશ થાય છે
(જેમ કે $\le$) અને કેટલાકમાં તે બાકાત રહે છે (જેમ કે $<$).
અહીં એક વર્ગીકરણ આપણને મદદ કરશે.
\end{explain}

\begin{defn}[પૂર્વક્રમ]
સ્વવાચક અને પરંપરિત બંને હોય તેવા સંબંધને \emph{પૂર્વક્રમ} કહે છે.
\end{defn}

\begin{defn}[આંશિક ક્રમ]
વિસંમિત પણ હોય તેવા પૂર્વક્રમને \emph{આંશિક ક્રમ} કહે છે.
\end{defn}

\begin{defn}[રેખીય ક્રમ]\ollabel{def:linearorder}
તુલનાયુક્ત પણ હોય તેવા આંશિક ક્રમને \emph{સંપૂર્ણ ક્રમ}
અથવા \emph{રેખીય ક્રમ} કહે છે.
\end{defn}

\begin{ex}
દરેક રેખીય ક્રમ આંશિક ક્રમ પણ છે, અને દરેક આંશિક ક્રમ
પૂર્વક્રમ પણ છે, પણ ઊલટાં વિધાનો સાચાં નથી.
$A$ પરનો સાર્વત્રિક સંબંધ પૂર્વક્રમ છે, કારણ કે તે સ્વવાચક
અને પરંપરિત છે. પણ, જો $A$માં એક કરતાં વધુ !!{element}
હોય, તો સાર્વત્રિક સંબંધ વિસંમિત નથી, અને તેથી આંશિક ક્રમ નથી.
\end{ex}

\begin{ex}
$\Bin^*$ પરનો \emph{કરતાં વધુ લાંબું નહીં} એવો સંબંધ
$\preccurlyeq$ લો: $x \preccurlyeq y$ તો અને તો જ
$\len{x} \le \len{y}$. આ પૂર્વક્રમ છે (સ્વવાચક અને પરંપરિત),
અને તુલનાયુક્ત પણ છે, પરંતુ આંશિક ક્રમ નથી,
કારણ કે તે વિસંમિત નથી. દાખલા તરીકે,
$01 \preccurlyeq 10$ અને $10 \preccurlyeq 01$,
પણ $01 \neq 10$.
\end{ex}

\begin{ex}
ગણોના ગણ પરનો સંબંધ $\subseteq$ એક મહત્વનો આંશિક ક્રમ છે.
સામાન્ય રીતે તે રેખીય ક્રમ નથી, કારણ કે જો $a \neq b$ હોય
અને આપણે $\Pow{\{a, b\}} = \{\emptyset, \{a\}, \{b\}, \{a,b\}\}$
લઈએ, તો જોવા મળે છે કે $\{a\} \nsubseteq \{b\}$ અને
$\{a\} \neq \{b\}$ અને $\{b\} \nsubseteq \{a\}$.
\end{ex}

\begin{ex}
\emph{નિઃશેષ વિભાજ્યતા}નો સંબંધ આપણને એવો આંશિક ક્રમ આપે છે
જે રેખીય ક્રમ નથી. પૂર્ણાંકો $n$ અને $m$ માટે આપણે $n \mid m$
લખીએ, જેનો અર્થ છે કે $n$ એ $m$ને (નિઃશેષ) ભાગે છે,
એટલે કે, જો અને ફક્ત જો કોઈક પૂર્ણાંક $k$ એવો હોય કે
$m = kn$. $\Nat$ પર આ આંશિક ક્રમ છે, પણ રેખીય ક્રમ નથી:
દાખલા તરીકે, $2 \nmid 3$ અને $3 \nmid 2$ પણ.
$\Int$ પરના સંબંધ તરીકે લઈએ તો વિભાજ્યતા માત્ર પૂર્વક્રમ છે,
કારણ કે તે વિસંમિત નથી: $1 \mid -1$ અને $-1 \mid 1$
પણ $1 \neq -1$.
\end{ex}

\begin{ex}
અનુક્રમોના ગણ $A^*$ પરનો \emph{વિસ્તરણ} સંબંધ આ છે:
$s \sqsubseteq s'$ તો અને તો જ $s = \emptyseq$ (ખાલી અનુક્રમ),
$s = s'$, અથવા $s = \tuple{s_1, \dots, s_n}$ અને
$s' = \tuple{s_1, \dots, s_n, s_{n+1}, \dots, s_m}$.
જો $s \sqsubseteq s'$, તો આપણે $s$ને $s'$નો
\emph{આરંભિક ખંડ} પણ કહીએ છીએ.
$A^*$ પરનો વિસ્તરણ સંબંધ આંશિક ક્રમ છે પણ રેખીય ક્રમ નથી,
દા.ત., જો $a \neq b$, તો $ab \not\sqsubseteq ba$ અને
$ba \not\sqsubseteq ab$.
\end{ex}

\begin{defn}[કડક ક્રમ]
\emph{કડક ક્રમ} એવો સંબંધ છે જે અસ્વવાચક, અસંમિત અને પરંપરિત હોય.
\end{defn}

\begin{defn}[કડક રેખીય ક્રમ]\ollabel{def:strictlinearorder}
તુલનાયુક્ત પણ હોય તેવા કડક ક્રમને \emph{કડક સંપૂર્ણ ક્રમ}
અથવા \emph{કડક રેખીય ક્રમ} કહે છે.
\end{defn}

\begin{ex}
$\le$ એ કડક રેખીય ક્રમ $<$ને અનુરૂપ રેખીય ક્રમ છે.
$\subseteq$ એ કડક ક્રમ $\subsetneq$ને અનુરૂપ આંશિક ક્રમ છે.
\end{ex}

$A$ પરના કોઈ પણ કડક ક્રમ $R$માં વિકર્ણ $\Id{A}$,
એટલે કે બધા જોડ $\tuple{x, x}$, ઉમેરીને તેને આંશિક ક્રમમાં
ફેરવી શકાય છે. (આને $R$નું \emph{સ્વવાચક સંવરણ} કહે છે.)
ઊલટું, આંશિક ક્રમમાંથી $\Id{A}$ દૂર કરીને કડક ક્રમ મેળવી શકાય છે.
હવે પછીનાં બે પરિણામો આ વાત ચોક્કસ કરે છે.

\begin{prop}\ollabel{prop:stricttopartial}
જો $R$ એ $A$ પરનો કડક ક્રમ હોય, તો $R^+ = R \cup \Id{A}$
આંશિક ક્રમ છે. વધુમાં, જો $R$ કડક રેખીય ક્રમ હોય, તો
$R^+$ રેખીય ક્રમ છે.
\end{prop}

\begin{proof}
ધારો કે $R$ કડક ક્રમ છે, એટલે કે $R \subseteq A^2$ અને
$R$ અસ્વવાચક, અસંમિત અને પરંપરિત છે.
$R^+ = R \cup \Id{A}$ લો. આપણે દર્શાવવું છે કે
$R^+$ સ્વવાચક, વિસંમિત અને પરંપરિત છે.

$R^+$ દેખીતી રીતે સ્વવાચક છે, કારણ કે દરેક $x \in A$
માટે $\tuple{x, x} \in \Id{A} \subseteq R^+$.

$R^+$ વિસંમિત છે તે દર્શાવવા, પરસ્પરવિરોધ મેળવવા માટે ધારો કે
$R^+xy$ અને $R^+yx$ પણ $x \neq y$.
$\tuple{x,y} \in R \cup \Id{A}$, પણ
$\tuple{x, y} \notin \Id{A}$ હોવાથી,
$\tuple{x, y} \in R$ હોવું જોઈએ, એટલે કે $Rxy$.
એ જ રીતે, $Ryx$. પણ આ $R$ અસંમિત છે એવી ધારણાનો વિરોધ કરે છે.

પરંપરિતતા સ્થાપવા માટે ધારો કે $R^+xy$ અને $R^+yz$.
જો $\tuple{x, y} \in R$ અને $\tuple{y,z} \in R$ બંને,
તો $R$ પરંપરિત હોવાથી $\tuple{x, z} \in R$.
નહિતર, કાં $\tuple{x, y} \in \Id{A}$, એટલે કે $x = y$,
અથવા $\tuple{y, z} \in \Id{A}$, એટલે કે $y = z$.
પ્રથમ કિસ્સામાં, ધારણા પ્રમાણે $R^+yz$ અને $x = y$,
તેથી $R^+xz$. બીજા કિસ્સામાં પણ એ જ રીતે.
બંને કિસ્સામાં $R^+xz$, તેથી $R^+$ પરંપરિત પણ છે.

“વધુમાં” વાક્યાંશ અંગે, ધારો કે $R$ તુલનાયુક્ત પણ છે.
તેથી દરેક $x \neq y$ માટે, કાં $Rxy$ અથવા $Ryx$,
એટલે કે, કાં $\tuple{x, y} \in R$ અથવા $\tuple{y, x} \in R$.
$R \subseteq R^+$ હોવાથી, $R^+$ માટે પણ આ સાચું રહે છે,
તેથી $R^+$ તુલનાયુક્ત પણ છે.
\end{proof}

\begin{prop}\ollabel{prop:partialtostrict}
જો $R$ એ $A$ પરનો આંશિક ક્રમ હોય, તો
$R^- = R \setminus \Id{A}$ કડક ક્રમ છે.
વધુમાં, જો $R$ રેખીય ક્રમ હોય, તો $R^-$ કડક રેખીય ક્રમ છે.
\end{prop}

\begin{proof}
આ સ્વાધ્યાય તરીકે છોડ્યું છે.
\end{proof}

\begin{prob}
\olref[sfr][rel][ord]{prop:partialtostrict}ની સાબિતી આપો.
\end{prob}

નીચેનું સરળ પરિણામ સ્થાપે છે કે કડક રેખીય ક્રમો,
ઘટકો દ્વારા નિર્ધારિત સમાનતાને મળતો આવતો ગુણધર્મ સંતોષે છે:

\begin{prop}\ollabel{prop:extensionality-strictlinearorders}
જો $<$ એ $A$ પરનો કડક રેખીય ક્રમ હોય, તો:
\[
  (\forall a, b \in A)((\forall x \in A)(x < a \liff x < b) \lif a = b).
\]
\end{prop}

\begin{proof}
ધારો કે $(\forall x \in A)(x < a \liff x < b)$.
જો $a < b$, તો $a < a$, જે $<$ અસ્વવાચક છે એ હકીકતનો
વિરોધ કરે છે; તેથી $a \nless b$. બરાબર એ જ રીતે,
$b \nless a$. તેથી $a = b$, કારણ કે $<$ તુલનાયુક્ત છે.
\end{proof}

\end{document}
